\chapter{General Medical Terms}

\medterm{Abdomen}-- The body region between the thorax and pelvis, containing the digestive organs, liver, gallbladder, pancreas, spleen, kidneys, and major blood vessels. Divided into nine regions (right/left hypochondriac, epigastric, umbilical, right/left lumbar, right/left inguinal) or four quadrants (RUQ, LUQ, RLQ, LLQ) for clinical description.

\medterm{Auscultation}-- The act of listening to sounds produced within the body, typically using a stethoscope. Essential for assessing heart sounds, lung sounds, bowel sounds, and vascular bruits. Normal findings: S1 and S2 heart sounds, vesicular breath sounds, normal bowel sounds. Abnormal: murmurs, crackles, wheezes, rubs, bruits.

\medterm{Biopsy}-- Removal of tissue for histological examination to diagnose disease. Types: needle biopsy (fine needle aspiration, core needle), excisional, incisional, punch, endoscopic. Sites: breast, lung, liver, lymph node, bone marrow. Complications: bleeding, infection, seeding (rare).

\medterm{Comorbidity}-- Presence of one or more additional conditions co-occurring with a primary disease. Common in chronic illnesses (diabetes with hypertension, COPD with heart failure). Affects treatment decisions, prognosis, and healthcare utilization.

\medterm{Contraindication}-- A condition or factor that renders a particular procedure, treatment, or medication inadvisable. Absolute contraindication: condition where intervention should never be performed. Relative contraindication: risks outweigh benefits but may be acceptable in specific circumstances. Examples: pregnancy (absolute for isotretinoin), severe renal impairment (relative for contrast CT).

\medterm{Diff diagnosis}-- The process of distinguishing a particular disease or condition from others that present with similar symptoms. Systematic approach: consider life-threatening conditions first, use pre-test probability, apply diagnostic criteria, order targeted tests.

\medterm{Edema}-- Swelling caused by excess fluid trapped in body tissues. Types: pitting (press leaves indentation), non-pitting (myxedema, lymphedema). Causes: heart failure, kidney disease, liver cirrhosis, venous insufficiency, lymphatic obstruction, medications (calcium channel blockers, NSAIDs). Assessment: degree, distribution, symmetry, associated symptoms.

\medterm{Examination}-- Systematic assessment of a patient's health status. Components: inspection, palpation, percussion, auscultation (IPPA). General survey: appearance, posture, gait, vital signs. Focused examination based on chief complaint.

\medterm{Homeostasis}-- The maintenance of stable internal conditions despite external changes. Key regulated variables: temperature, pH, glucose, electrolytes, blood pressure, fluid balance. Mechanisms: negative feedback loops (most common), positive feedback (rare, e.g., labor, blood clotting). Organs involved: hypothalamus, kidneys, lungs, liver, endocrine glands.

\medterm{Incidence}-- Number of new cases of a disease in a population during a specified time period. Measures risk of developing disease. Incidence rate = new cases / population at risk × time unit.

\medterm{Inflammation}-- The body's protective response to injury or infection. Cardinal signs: redness (rubor), heat (calor), swelling (tumor), pain (dolor), loss of function (functio laesa). Mediators: histamine, prostaglandins, cytokines, complement. Acute: neutrophil-predominant. Chronic: lymphocyte/macrophage-predominant.

\medterm{Informed Consent}-- The process by which a patient voluntarily confirms understanding of a proposed treatment or procedure, including risks, benefits, alternatives, and consequences of non-treatment. Elements: disclosure, comprehension, voluntariness, competence, consent. Documented on consent form. Exceptions: emergencies, waived consent (minimal risk research), therapeutic privilege.

\medterm{Iatrogenic}-- An illness or condition caused by medical examination or treatment. Examples: hospital-acquired infections, drug side effects, surgical complications, radiation-induced damage. Prevention: proper technique, informed consent, monitoring, reporting systems.

\medterm{Jargon}-- Specialized language used by a particular group or profession. Medical jargon can hinder communication with patients. Examples: "stat" (immediately), "NPO" (nothing by mouth), "DX" (diagnosis). Modern trend toward plain language in patient communication.

\medterm{Kussmaul Respirations}-- Deep, labored breathing pattern associated with diabetic ketoacidosis. Compensatory hyperventilation to blow off CO2 and correct metabolic acidosis. Name: Adolf Kussmaul, German physician.

\medterm{Lateral}-- Positioned away from the midline of the body. Opposite of medial. Example: the thumb is lateral to the little finger in anatomical position.

\medterm{Mediate}-- To transmit or convey between two points. In medicine, mediation often refers to the mechanism by which a cause produces an effect. Mediated immunity involves immune cells; direct contact does not.

\medterm{Nadir}-- The lowest point, particularly referring to the lowest white blood cell count after chemotherapy. Bone marrow suppression typically reaches nadir 7-14 days post-treatment. Important for infection risk assessment.

\medterm{Oropharynx}-- The posterior portion of the mouth leading to the esophagus, containing the tonsils and base of tongue. Area for swallowing and speech. Examination with tongue depressor and light source. Pathology: pharyngitis, tonsillitis, malignancy.

\medterm{Palpation}-- Examination by touch to detect abnormalities in size, shape, consistency, and tenderness of organs or tissues. Uses fingers, palm, or heel of hand. Light vs. deep palpation. Bimanual palpation uses two hands.

\medterm{Quadrant}-- One of four sections of the abdomen (RUQ, LUQ, RLQ, LLQ) used for clinical description. RUQ: liver, gallbladder, duodenum, right kidney. LUQ: stomach, spleen, left kidney, pancreas tail. RLQ: appendix, cecum, right ovary. LLQ: sigmoid colon, left ovary.

\medterm{Rx}-- Prescription or treatment. Symbol derives from Latin "requirere" (to take). Also used as abbreviation for remedy. Rx symbol: Rx or &x;.

\medterm{Sensitivity}-- Proportion of true positives correctly identified by a test. High sensitivity: few false negatives. Useful for screening. Formula: true positives / (true positives + false negatives).

\medterm{Tincture}-- Alcohol extract of plant or chemical substance. Common concentration: 10-40% alcohol. Examples: tincture of iodine (antiseptic), tincture of benzoin (skin protector). Also used for liquid medications.

\medterm{Ulnar}-- Relating to the ulna bone of the forearm or the ulnar side (medial side) of the hand. Ulnar nerve compression causes paresthesia in medial one and a half fingers.

\medterm{Vaccination}-- Administration of a vaccine to stimulate immune response and provide protection against infectious disease. Types: live attenuated, inactivated, subunit, mRNA, viral vector. Schedule varies by country. Herd immunity achieved when sufficient proportion immune.

\medterm{Wedge resection}-- Surgical removal of a wedge-shaped piece of tissue, typically lung or breast. Used for diagnosis or treatment of localized lesions. Less extensive than lobectomy.

\medterm{Xerosis}-- Abnormal dryness of skin or mucous membranes. Causes: aging, environmental factors, medications, systemic disease (hypothyroidism, diabetes), nutritional deficiencies. Treatment: emollients, humidification, treat underlying cause.

\medterm{Yield}-- The proportion of positive test results among those tested. In clinical trials, yield refers to the proportion of patients benefiting from intervention. Diagnostic yield: percentage of tests providing diagnostic information.

\medterm{Zoster}-- See herpes zoster (shingles).

\medterm{Auscultatory Gap}-- A transient silence between the first and second sounds of blood pressure measurement, sometimes heard in hypertensive patients. Can lead to underestimation of systolic pressure if not recognized.

\medterm{Barium Swallow}-- Radiographic examination of the esophagus using barium sulfate contrast. Evaluates swallowing function, motility, strictures, masses. Modified barium swallow assesses aspiration risk.

\medterm{Capillary Refill Time}-- Time for color to return to blanched nail bed or skin after pressure. Normal: <2 seconds. Prolonged in shock, dehydration, peripheral vasoconstriction.

\medterm{Diuresis}-- Increased urine production. Can be physiological (fluid overload) or pharmacological (diuretics). Osmotic diuresis (hyperglycemia), water diuresis (DI), saline diuresis (hypertonic saline).

\medterm{Echo}-- See echocardiography.

\medterm{Fistula}-- Abnormal connection between two epithelial surfaces. Examples: arteriovenous fistula, enterocutaneous fistula, rectovaginal fistula. Can be congenital or acquired (surgery, infection, trauma, malignancy).

\medterm{Gait}-- Pattern of walking. Normal gait involves coordinated movement of legs, arms, trunk. Abnormal gait patterns: antalgic (pain), hemiplegic (stroke), parkinsonian (shuffling), ataxic (cerebellar), steppage (foot drop).

\medterm{Hemostasis}-- The process of stopping bleeding. Primary: platelet plug formation. Secondary: coagulation cascade. Tertiary: vasoconstriction and fibrinolysis regulation.

\medterm{Inspection}-- Visual examination of the patient. Part of physical assessment. Observes symmetry, color, lesions, movement, posture.

\medterm{Joint Effusion}-- Excess fluid within a joint space. Common sites: knee (most accessible), ankle, wrist, shoulder. Causes: trauma, infection, inflammatory arthritis, crystal arthropathy. Detected by ballottement or bulge sign.

\medterm{Kernig Sign}-- Sign of meningeal irritation. Patient supine, hip and knee flexed 90 degrees; inability to extend knee indicates positive sign. Seen in meningitis, subarachnoid hemorrhage.

\medterm{Leukocytosis}-- Elevated white blood cell count (>11,000/\mu L). Causes: infection, inflammation, stress, medications (corticosteroids), malignancy (leukemia). Neutrophilia, lymphocytosis, eosinophilia, monocytosis have different implications.

\medterm{Murmur}-- Heart sound caused by turbulent blood flow. Grades I-VI. Systolic murmurs more common; diastolic murmurs always pathological. Evaluation: timing, location, radiation, intensity, quality, position changes.

\medterm{Nodule}-- Small, palpable mass <2 cm. Can be solid or cystic. Examples: thyroid nodule, rheumatoid nodule, pulmonary nodule. Evaluation depends on location.

\medterm{Oliguria}-- Decreased urine output (<400 mL/day in adults). Causes: prerenal (dehydration, heart failure), renal (acute kidney injury), postrenal (obstruction). Medical emergency if anuric.

\medterm{Percussion}-- Tapping body surface to assess underlying structures. Sounds: dull (solid organ, fluid), resonant (lung), tympanic (gas-filled). Used to assess organ size, presence of fluid, air.

\medterm{Quiescent}-- Inactive or dormant state of disease. Example: quiescent ulcerative colitis (no symptoms but histologic inflammation may persist).

\medterm{Referred Pain}-- Pain perceived at location other than site of pathology. Examples: myocardial infarction (left arm/jaw), gallbladder (right shoulder), diaphragmatic irritation (shoulder).

\medterm{Saturate}-- To absorb or dissolve as much as possible. Oxygen saturation (SpO2) measures percentage of hemoglobin binding sites occupied by oxygen.

\medterm{Tachypnea}-- Rapid breathing (>20 breaths/min in adults). Causes: fever, pain, anxiety, pulmonary disease, metabolic acidosis. Assessment of rate, depth, pattern essential.

\medterm{Urination}-- See micturition.

\medterm{Vital Signs}-- Key physiological measurements: temperature, pulse, respiration rate, blood pressure, oxygen saturation, pain. Normal ranges: T 36.5-37.5°C, HR 60-100, RR 12-20, BP <120/80, SpO2 >95%.

\medterm{Wheal}-- Raised, itchy swelling of skin (urticaria/hives). Caused by mast cell degranulation and histamine release. Examples: insect bite, allergic reaction, dermatographism.

\medterm{X-ray}-- See radiography.

\medterm{Yield Stress}-- In biomechanics, the stress at which a material begins to deform plastically. Relevant to implant design and bone mechanics.

\medterm{Zollinger-Ellison Syndrome}-- See gastroenterology.

\medterm{Abscess}-- See emergency_medicine.

\medterm{Biopsy}-- See general.

\medterm{Comorbidity}-- See general.

\medterm{Diuresis}-- See general.

\medterm{Edema}-- See general.

\medterm{Fever}-- Elevated body temperature (>38°C or 100.4°F). Workup depends on duration, height, and patient factors. Common causes: infections (bacterial, viral), autoimmune disorders, malignancies, drug reactions. Management: antipyretics (acetaminophen, ibuprofen), treat underlying cause. Febrile seizures: generally benign, prevent recurrence with fever control.

\medterm{Gait}-- See general.

\medterm{Hemostasis}-- See general.

\medterm{Infection}-- See infectious_disease.

\medterm{Joint Effusion}-- See general.

\medterm{Kussmaul Respirations}-- See general.

\medterm{Leukocytosis}-- See general.

\medterm{Murmur}-- See general.

\medterm{Nodule}-- See general.

\medterm{Oliguria}-- See general.

\medterm{Percussion}-- See general.

\medterm{Referred Pain}-- See general.

\medterm{Saturate}-- See general.

\medterm{Tachypnea}-- See general.

\medterm{Urination}-- See general.

\medterm{Vital Signs}-- See general.

\medterm{Wheal}-- See general.

\medterm{Xerosis}-- See general.

\medterm{Yield}-- See general.

\medterm{Zoster}-- See infectious_disease.
